\chapter{Introduction}

\paragraph{Privacy and power.} Resources are not equally distributed throughout
the world. There are infinitely many consequences of this. All of them
influence our actions in some way, but not all of them can be influenced in any
way by cryptography.

Encryption allows people to communicate by exchanging messages that are, to
eavesdroppers, indistinguishable from random. This allows parties with a shared
key to communicate privately, regardless of their personal power, wealth,
affiliations, or any other characteristics.
% A shifting sandbank in eternity,
The desire to conceal written information has been considered for at least the
last four thousand years, reaching back to the obfuscated hieroglyphs in the
tomb of Khnumhotep II.  Modern cryptography has enabled encrypted
communications that even truly powerful entities cannot read.

The communicating parties themselves are not assumed to be as powerful as the
entities trying to break the system.  Encryption demands relatively little
computational power of the users while keeping the adversary's computational
task infeasible.  This asymmetry of necessary computational power between users
and attacking entities is well noted, and lies at the foundation of modern
cryptography itself.

Cryptography today deals in more varied technologies than private
communication.  Information flows, and the question shifts to reflect that.
Once the contents of a communication are protected, for example, a resourceful
adversary may instead attempt to determine the identities of the communicating
parties, their relationships, or patterns in their behaviour. Protecting that
then brings into focus the next piece of information that could be observed and
used. Every answer reveals another question.

Over the last fifty years, privacy enhancing technologies have been designed to
answer exactly these broader privacy problems. The surface of information being
protected has expanded, and so has the number of entities capable of observing
some part of the system.  Service providers, issuers, validators, network
operators and public authorities may all learn sensitive information while
performing otherwise legitimate functions.

These powerful entities are consequently often modelled as adversaries.
Protocol designers have then, understandably, been active in designing systems
that take power away from powerful, centralised entities. Anything that such an
entity learns is a step away from the perfect world. The ideal world.

\paragraph{Centralisation, responsibility and trust.} In attempting to remove
trust from centralised institutions, protocols may instead transfer
computation, storage and operational responsibility to individual users. This
can improve some properties of a system, but centralisation is not a single
property and decentralisation is not a complete description of a protocol's
design.

Computation, state, authority and information may be distributed differently.
A system may centralise computation while distributing authority between
several institutions, or centralise storage while using cryptography to prevent
the storage provider from learning the stored information. The distribution of
these responsibilities does not by itself determine the trust assumptions of
the system. An institution may be relied upon to provide availability, maintain
state or execute a prescribed computation without being trusted with users'
private information.

Privacy-preserving systems are often discussed as if the choice is between
decentralised privacy and centralised surveillance. This thesis challenges that
binary. Powerful institutions can perform useful protocol functions while
cryptography limits both the information available to them and the capabilities
they receive. Allowing institutions to perform this work may reduce the
computation, storage, time and specialist knowledge required of individual
users.

Modern privacy-preserving systems already involve service providers, issuers,
auditors, validators and other accountable entities with operational
responsibilities. Privacy Pass, anonymous credential systems and privacy
services such as Apple's Private Relay use institutional infrastructure while
limiting the information available to the institutions
involved~\cite{PoPETS:DGSTV18,privaterelay2021}. This thesis explores how
modelling such roles explicitly can enable practical and efficient
privacy-preserving systems.

\paragraph{Financial privacy.} Cryptographic protocols are applied in vastly
different spaces. There are extremely strong arguments against governments
being able to read the communications of their citizens, and companies the
messages of their customers.  A private communication system can aim to ensure
that no party other than the intended recipient learns the contents of a
message.

But in a financial system, there is something slightly different involved.
Real financial systems have legitimate accountability requirements.

They must establish the validity of transactions, maintain a consistent shared
state, control issuance, prevent or detect double spending, and, in many
settings, support some form of accountability.  This thesis takes the existence
of financial auditing requirements as a system constraint and asks how narrowly
the corresponding capabilities can be defined.  What information should an
auditor be able to obtain? Under which circumstances can they obtain it?
Should access be limited to a transaction, a user, a certain period of time, a
single prescribed purpose, or be ongoing?

There are often objections to constructing protocols with accountability at
all, as it seems analogous to a backdoor in the system. And there is the
argument that once a capability exists, it will be used everywhere possible.
This is important to consider.  Financial systems are inherently responsible
for movement of money.  This instantly distinguishes the system from
communication alone.  Refusing to design auditing protocols does not remove the
accountability requirements already present in financial systems. It may
instead result in the auditing goals being met by systems offering much weaker
privacy to all users.

This thesis starts with the presence of these operational roles and asks how
cryptography can work with them. An institution need not be modelled in the
same way for every property of a system. For availability and correct
operation, we assume that it performs its prescribed role. For privacy, the
same institution may deviate arbitrarily from the protocol, retain and combine
the information available to it, and attempt to learn anything it can. Privacy
is assessed against the strongest information it can derive in this setting.
The institution may therefore be relied upon to perform an operational function
without being trusted with users' private information.

\paragraph{The thesis.} This thesis argues that practical privacy is achieved
not by assuming away existing infrastructure, but by modelling its participants
and responsibilities accurately, and cryptographically constraining what each
participant can learn and do. In particular, an institution may perform
computation, maintain state, issue credentials, validate transactions or
provide availability without being trusted with users' private information. We
examine this principle through several systems.

First, we consider aggregate anonymous credentials. Existing credential models
do not always reflect the way that authority is divided between institutions in
practice. A user may receive attributes from several independent issuers and
later need to demonstrate statements concerning those attributes without
revealing either the attributes themselves or unnecessary information about the
issuers. We construct an efficient multi-issuer aggregate anonymous credential
scheme in which independently issued credentials can be combined and presented
through a single proof. Anonymous credential systems already implicitly include
a registration authority. We use this existing role to bind credentials issued
to the same user, preventing credentials belonging to different users from
being aggregated without requiring expensive proofs during presentation.

We next consider accountability in privacy-preserving transaction systems.
Financial systems may require both the ability to investigate a particular
transaction and the ability, under prescribed circumstances, to identify the
transactions associated with a particular user. These capabilities are related
but distinct, and granting one need not grant the other. We develop a
transaction identification mechanism that supports concurrent transaction
creation without requiring the auditor to maintain persistent per-user state.
Tracing authority can be restricted to a specified epoch, limiting both the
scope and duration of the information disclosed to an auditor.

Third, we consider transaction authorisation in permissioned distributed
systems. In these systems, transactions are approved by known endorsing parties
rather than anonymous participants. We construct threshold structure-preserving
signature schemes that allow the required endorsers to jointly authorise
transactions without relying on a single signing authority.  The schemes
require no interaction during the online signing phase, making them suitable
for efficient transaction authorisation. This treats the existing institutional
structure as a resource to be modelled rather than an obstacle to be removed.

Finally, we study privacy-preserving offline payments. Offline operation
improves availability and accessibility, but changes how double spending must
be handled. We develop a privacy-preserving offline payment scheme that divides
computation between a mobile device and a secure element. Payments made by
honest users remain unlinkable, while information identifying the double
spender can be recovered following double spending.

Together, these contributions support three claims:

\begin{enumerate}

  \item Centralisation is multidimensional. Computation, state, authority, and
    knowledge can each be centralised or decentralised independently, with
    substantially different consequences for the resulting system.

  \item Protocol models should reflect the roles that exist in practice.
    Issuers, auditors, validators and service providers can be represented
    explicitly rather than removed in pursuit of decentralisation.

  \item Privacy does not automatically follow from decentralisation, nor does
    it automatically disappear when corporations or public institutions
    participate in a protocol. What matters is the information made available
    to each participant.

\end{enumerate}

The remainder of this thesis is organised as follows.
\Fref{chap:preliminaries} introduces the notation, security definitions and
cryptographic building blocks used throughout. \Fref{chap:agg-creds} considers
privacy-preserving identity systems and presents our aggregate anonymous
credential construction. \Fref{chap:txid} presents our transaction
identification mechanism.  \Fref{chap:tsps} presents our threshold
structure-preserving signature protocol, and \Fref{chap:offline-payments}
presents our offline payment mechanism. 
Finally, \Fref{chap:conclusion} summarises the work and discusses open
questions.
